\section{Discussion}
\label{sec:discussion}

\subsection{Choosing $\alpha$}

The stickiness parameter $\alpha$ traces a Pareto frontier between two objectives:
minimising boundary length (compactness, $\alpha = 0$) and minimising county
splits (subdivision preservation, $\alpha \to \infty$).
Figure~\ref{fig:frontier} (forthcoming) plots this frontier for representative
states across a range $\alpha \in [0, 50]$.

For practical redistricting, we suggest $\alpha = 5$ as a default.
At this value:
\begin{itemize}
  \item County splits are halved relative to MEC ($-45\%$)
  \item Edge cut increases by a factor of 2.5 --- still well within the range
    of plans courts have accepted as compact
  \item The parameter is interpretable: an intra-county edge costs $1 + 5 = 6$
    times its geographic boundary length, while a cross-county boundary costs 1
  \item No partisan bias is introduced
\end{itemize}

Statutory language varies: California says ``respect the geographic integrity''
of subdivisions; North Carolina requires using whole counties wherever possible;
Florida prohibits splitting municipalities.
A redistricting body could choose $\alpha$ based on how strictly its statute
mandates preservation.

\subsection{Why County Splits Are Not Zero at $\alpha = 5$}

Even with $\alpha = 5$, some county splits remain.
This is a consequence of the population balance constraint ($\pm 0.5\%$ for
congressional districts).
When a county has population well above the ideal district size (e.g.,
Los Angeles County, population 10M, contains about 15 congressional districts),
it must be split regardless of $\alpha$.
The floor on unavoidable splits is approximately:
\[
  \text{min\_splits}(\text{state}) \approx \sum_c \max(0, \lceil P_c / P^* \rceil - 1)
\]
where $P_c$ is county population and $P^*$ is the ideal district population.
At $\alpha = 5$, many states approach this theoretical minimum.

\subsection{Partisan Neutrality}

The absence of systematic partisan bias at $\alpha = 5$ is not obvious.
Counties in the United States are not politically neutral geographic units:
urban counties tend to vote Democratic, rural counties Republican.
One might expect that respecting county lines would ``pack'' urban Democratic
voters into fewer districts.

The data shows otherwise.
The reason is that METIS with $\alpha > 0$ does not simply consolidate counties;
it makes county-crossing cuts more expensive while still minimising total cut
weight subject to balance.
The algorithm finds the most compact county-respecting partition, which does
not systematically favour either party because county shapes and populations
are not themselves gerrymandered.

This is the key legal point: $\alpha$ operationalises a traditional criterion
that courts have consistently upheld as a legitimate redistricting objective,
and it does so without introducing partisan information.

\subsection{Extension to Finer Levels}

The framework extends to county subdivisions (townships, MCDs), incorporated
places (cities, towns), and voting tabulation districts (precincts).
Each level requires an additional TIGER shapefile download and a spatial join
to assign each census tract to its containing jurisdiction.
We defer this to future work; the county level alone (derivable from existing
data) produces the majority of the reduction in splits.

\subsection{Limitations}

The current evaluation uses a single seed ($= 1$) for each state-$\alpha$ pair.
As shown in B.7~\citep{DL2026a}, single-seed MEC results may not represent the
converged minimum.
A full convergence analysis would require running 100--500 seeds per state per
$\alpha$ value.
We regard the present results as a first-pass characterisation of the Pareto
frontier shape; convergence is left for the full analysis.

Additionally, the partisan outcomes depend on the specific plan found (which
depends on seed), so the partisan table should be interpreted as indicative
rather than definitive.
